\documentclass{article}
\usepackage{amsmath, amssymb, amsthm}
\usepackage{hyperref}
\usepackage{geometry}
\geometry{margin=1in}

\title{Erd\H{o}s Problem \#132}
\date{Last updated: 2025-08-31}
\author{Source: erdosproblems.com}

\begin{document}
\maketitle

\noindent\textbf{Status:} OPEN \quad \textbf{Prize: \$100}

\noindent\textbf{Tags:} distances

\medskip
\noindent\textbf{Problem Statement:}

Let $A\subset \mathbb{R}^2$ be a set of $n$ points. Must there be two distances which occur at least once but between at most $n$ pairs of points? Must the number of such distances $\to \infty$ as $n\to \infty$?

\bigskip
\noindent\textbf{Remarks:}

Asked by Erd\H{o}s and Pach. Hopf and Pannowitz \cite{HoPa34} proved that the largest distance between points of $A$ can occur at most $n$ times, but it is unknown whether a second such distance must occur.

It may be true that there are at least $n^{1-o(1)}$ many such distances. In \cite{Er97e} Erd\H{o}s offers \$100 for 'any nontrivial result'.

Erd\H{o}s \cite{Er84c} believed that for $n\geq 5$ there must always exist at least two such distances. This is false for $n=4$, as witnessed by two equilateral triangles of the same side-length glued together. Erd\H{o}s and Fishburn \cite{ErFi95} proved this is true for $n=5$ and $n=6$.

Clemen, Dumitrescu, and Liu \cite{CDL25} have proved that there always at least two such distances if $A$ is in convex position (that is, no point lies inside the convex hull of the others). They also prove it is true if the set $A$ is 'not too convex', in a specific technical sense.

See also [223], [756], and [957].

\bigskip
\noindent\textbf{References:}

[CDL25] F. Clemen, A. Dumitrescu, and D. Liu, On multiplicities of interpoint distances. arXiv:2505.04283 (2025).
 
 [Er84c] Erd\H{o}s, Paul, Some old and new problems in combinatorial geometry. Convexity and graph theory (Jerusalem, 1981) (1984), 129-136.
 
 [Er97e] Erd\H{o}s, Paul, Some of my favourite unsolved problems. Math. Japon. (1997), 527-537.
 
 [ErFi95] Erd\H{o}s, Paul and Fishburn, Peter C., Multiplicities of interpoint distances in finite planar sets. Discrete Appl. Math. (1995), 141--147.
 
 [HoPa34] Hopf, H. and Pannwitz, E., Aufgabe 167. Jber. Deutsch. Math. Verein. (1934), 114.

\bigskip
\noindent\small{Source: \url{https://www.erdosproblems.com/132}}
\end{document}
